\section{Finite-circle Bethe states and the transfer norm}\label{sec:bethe}

We now complete the spectral calculation required by
Proposition~\ref{prop:brownian-parity}.  Lemma~\ref{lem:brownian-ground}
reduces the norm of $T_{n_p}$ to its eigenvalue on the balanced ground state
of $H_2$.  The remaining tasks are to construct that state on the actual
circle, to identify it without a completeness assertion for periodic Bethe
states, and to evaluate the transfer eigenvalue including its winding terms.
The construction uses two scales.  At each \emph{fixed} particle number,
a line spectral resolution will identify the possible strong-coupling
limits on a circle.  We then construct finite-circle roots, prove that
their wavefunction has a strictly positive limiting shape,
and continue that identification to the required circumference.
The root and transfer estimates, in contrast to this identification
argument, must be uniform as particle number grows with circumference.

Write
\[
 n_p=2M+\ell,\qquad M\in\mathbb Z_{\geq0},\qquad \ell\in\{0,1\}.
\]
Unless stated otherwise, the estimates involving $b$ below have the following
meaning.  Fix $D>0$, a constant $C_{\mathrm{band}}<\infty$, and positive lower
and upper constants in
\[
 W\asymp e^{\pi b}b^{-1/2}(\log b)^D,
 \qquad |n_p/W-b^2|\leq C_{\mathrm{band}}b.
\]
Then $b$ is sufficiently large in terms of these constants.  The constants
in uniform estimates may depend on these fixed data, but not on $n_p$, $W$,
or the larger circumference $\mathcal W\geq W$ used in the continuation.
The notation $\operatorname{poly}(b)$ denotes a fixed polynomial bound,
whose degree may increase after finitely many differentiations.  Bounds
with a subscript $n_p$ are instead used only with $n_p$ fixed.

\subsection{Contact scattering and the finite-dimensional spin calculation}

The coordinate and nested Bethe-ansatz methods for one-dimensional
contact interactions originate in the work of Lieb and Liniger,
Gaudin, and Yang~\cite{LiebLiniger1963,Gaudin1967,Yang1967}.
We give the scattering and auxiliary-spin calculations in the present
normalization, since both their signs and their fusion factors enter
the finite-circle proof.

In the ordered coordinate chamber $X_1<\cdots<X_{n_p}$, an eigenfunction
with momenta $k_1,\ldots,k_{n_p}$ is a sum of plane waves, one for each
assignment of momenta to the slots.  Its coefficients lie in
$(\mathbb C^2\otimes\mathbb C^2)^{\otimes n_p}$.  The permutations $P^L$
and $P^R$ below exchange, respectively, the left and right binary indices
of two neighboring slots.  If their assigned momenta have difference $x$,
exchanging those momenta changes the coefficient by
\begin{equation}\label{eq:bethe-R}
 \check R(x)=P^LP^R
       \frac{(x-icP^L)(x-icP^R)}{x^2+c^2}.
\end{equation}
Indeed, in the color subspace symmetric under $P^LP^R$, the contact condition
for \eqref{eq:B2} is
\[
 (\partial_{a+1}-\partial_a)\psi=cP^L\psi
       \quad\hbox{at }X_{a+1}=X_a;
\]
the complementary color subspace has Dirichlet boundary condition.  Substituting
the two plane waves at this face gives \eqref{eq:bethe-R}.  The inverse
relation and the braid relation follow by multiplying the two single-index
relations for $x-icP$.  Each single-index identity follows by expanding and
using the permutation relations.  Consequently the coefficient obtained
from an initial coefficient does not depend on the chosen sequence of
neighboring exchanges.

We use momentum units in which $c=2$ and denote the circumference in those
units by $\mathcal W$.  A boundary twist will remove apparent singularities
in intermediate algebraic formulas.  It is
\[
 e^{i\zeta}g\otimes g,
 \qquad g=\operatorname{diag}(e^{i\chi},e^{-i\chi}),
\]
with $\zeta,\chi$ near zero.  Introduce $M$ auxiliary variables $a_r$ and the
polynomial $Q_0(u)=\prod_{r=1}^M(u-a_r)$.  We use the same auxiliary variables
for the two binary indices.

\begin{lemma}[Algebraic eigenvectors]\label{lem:bethe-algebra}
Suppose the momenta and auxiliary variables are nonsingular solutions of
\begin{equation}\label{eq:B8}
 \begin{aligned}
 e^{ik_j\mathcal W}
   &=e^{i\zeta+2i\chi}
       \prod_{t\ne j}\frac{k_j-k_t+2i}{k_j-k_t-2i}
       \left(\frac{Q_0(k_j-i)}{Q_0(k_j+i)}\right)^2,
       &&1\leq j\leq n_p,\\
 e^{-2i\chi}\prod_{j=1}^{n_p}
       \frac{a_r-k_j-i}{a_r-k_j+i}
   &=\prod_{s\ne r}\frac{a_r-a_s-2i}{a_r-a_s+2i},
       &&1\leq r\leq M.
 \end{aligned}
\end{equation}
Here nonsingular means that the displayed rational functions and the
coefficient formulas in the proof are defined by ordinary substitution.
For each binary index, apply the upper-right entries $B(a_1),\ldots,B(a_M)$
of the spin-$\frac12$ monodromy defined below to the all-up vector.  The
tensor product of these two vectors, propagated by \eqref{eq:bethe-R},
gives a periodic twisted eigenfunction of $H_2$ whenever it is nonzero.
Its energy is $\sum_j k_j^2$.

More generally, let $S_J$ be the spin-$J$ angular momentum matrices,
$g_J$ the spin-$J$ representation of $g$, $\sigma$ the Pauli matrices, and
\[
 \mathcal L_J(x)=x-2iS_J\cdot\sigma,
 \qquad F(u)=\prod_{j=1}^{n_p}(u-k_j).
\]
In the product of $\mathcal L_J(u-k_j)$, put later physical sites on the
left.  Its auxiliary trace with twist $g_J^{-1}$ has eigenvalue
\begin{equation}\label{eq:B9}
 t_J(u)=\sum_{m=-J}^{J}e^{-2im\chi}F(u-2im)
 \frac{Q_0(u-i(2J+1))Q_0(u+i(2J+1))}
      {Q_0(u-i(2m+1))Q_0(u-i(2m-1))}.
\end{equation}
All these identities continue meromorphically in the variables.  A
nonzero continuation limit of the plane eigenfunction is again an
eigenfunction at the limiting parameters.
\end{lemma}

\begin{proof}
For $J=\frac12$, the local matrix is
$\mathcal L_{1/2}(x)=(x+i)-2iP$.  Write its monodromy entries as $A,B,C,D$.
The relation with intertwiner $u-v-2iP$ shows first that the $B$ entries
commute.  It also gives the following coefficients when a diagonal entry
is moved to the right past $B(a)$: the terms that retain its spectral
variable $u$ have factors
\[
 \frac{u-a+2i}{u-a}\quad\hbox{for }A(u),
 \qquad
 \frac{u-a-2i}{u-a}\quad\hbox{for }D(u).
\]
The terms that replace this variable by $a$ have the opposite pole
coefficients and contain $B(u)A(a)$ or $B(u)D(a)$.  The all-up eigenvalues
of $A(u)$ and $D(u)$ are $F(u-i)$ and $F(u+i)$.

For several $B$ entries, first commute a selected $B(a_r)$ past the other
$B$'s.  The coefficient of the term in which $B(a_r)$ is replaced by
$B(u)$ is its pole coefficient times the keep factors for every other
auxiliary variable, evaluated at $a_r$.  This also follows directly by
induction and partial fractions.  The two such coefficients, from $A$
and $D$ in the twisted trace, cancel precisely by the second equation
of \eqref{eq:B8}.  The terms with no replacement give \eqref{eq:B9}
for $J=\frac12$.

For the higher auxiliary spins, multiply
\[
 \mathcal L_{1/2}(x+2iJ)\mathcal L_J(x-i).
\]
Relative to the two total auxiliary spins $J+\frac12$ and $J-\frac12$,
this product is triangular with diagonal blocks
\[
 (x+i(2J-1))\mathcal L_{J+1/2}(x),
 \qquad
 (x+i(2J+1))\mathcal L_{J-1/2}(x-2i).
\]
To verify the blocks, let $A$ and $B$ be the two auxiliary angular momenta.
The Pauli identity and $[A\cdot B,A]=iA\times B$ reduce the expansion to
the two total-spin projections; the projection of $A$ on these blocks is
$\pm(A+B)/(2J+1)$.  The same calculation proves invariance of the plus
block.  Taking twisted traces yields a recurrence in $J$.  Substitution
of \eqref{eq:B9} verifies the recurrence and its $J=0,\frac12$ initial
values, proving the formula for all $J$.

It remains to impose the boundary condition.  Move a momentum initially
in the first slot once around the circle.  The spin-$\frac12$ trace is
tested at $k_j-i$, where the local factor at that slot is a permutation
times $-2i$.  In each index the resulting transport is the twisted trace
divided by $-2i$; the joint scalar denominator from the other slots is
\[
 \prod_{t\ne j}\bigl((k_j-k_t)^2+4\bigr).
\]
Together with the plane-wave phase this is exactly the first equation
of \eqref{eq:B8}.  Braid consistency allows any momentum to be moved into
the first slot, so every boundary condition follows.  The contact
conditions were already imposed by \eqref{eq:bethe-R}.  Finally the
coefficients are rational functions times entire plane factors, which
proves the asserted meromorphic continuation.  At a nonzero limit the
contact and periodic boundary conditions persist, as does the eigenvalue
equation.
\end{proof}

The algebra constructs candidate eigenstates.  It does not identify their
position in the spectrum.  We next obtain the spectral information needed
for that identification on the line, where contour integration also
explains why only single particles and singlet pairs appear.

\subsection{A spectral resolution on the line}
\label{subsec:bethe-line}

The ground-state identification will start at strong coupling, where each
singlet pair has binding energy \(-c^2/2\).  We need to know that every
state with energy \(-Mc^2/2+O(1)\) is concentrated on \(M\) such pairs,
with at most one single particle.  We prove this first on the line by a
spectral resolution, and then pass to a fixed circle by localization.
Throughout this argument the particle number \(n_p=2M+\ell\) is fixed,
with \(\ell\in\{0,1\}\), and the coupling \(c>0\) is not scaled to \(2\).

Write
\[
 \mathcal C_{n_p}=\{X\in\mathbb R^{n_p}:X_1<\cdots<X_{n_p}\},
 \qquad
 \mathcal V_{n_p}=(\mathbb C^2\otimes\mathbb C^2)^{\otimes n_p}.
\]
We realize the line Hamiltonian \(H_c\) on
\(L^2(\mathcal C_{n_p};\mathcal V_{n_p})\), with the contact conditions
already derived from \eqref{eq:B2}.  Lebesgue measure on the ordered
chamber is used throughout.  For two slots \(a,b\), let
\(s_{ab}=2^{-1/2}(e_1\otimes e_2-e_2\otimes e_1)\) be the normalized
singlet in one binary index, and let
\(s_{ab}^{L}\otimes s_{ab}^{R}\) be the double singlet.

For integers \(\ell',m'\geq0\) satisfying \(\ell'+2m'=n_p\), choose
\(\ell'\) real elementary momenta \(k_1,\ldots,k_{\ell'}\) and \(m'\)
real pair centers \(u_1,\ldots,u_{m'}\).  The two momenta belonging to
center \(u_a\) are
\[
                 u_a-\frac{ic}{2},\qquad u_a+\frac{ic}{2}.
\]
Initially put each pair in adjacent slots, with its lower momentum
first, and insert its double singlet.  The elementary slots retain
arbitrary colors.  The resulting plane-wave sum, denoted
\(E_X^{\ell'm'}(k,u)\), is an operator
\(\mathcal V_{\ell'}\to\mathcal V_{n_p}\).  Its coefficients are obtained
with the interchange matrix \(\check R\); the sum runs over all
shuffles preserving the order of the two members of every pair.
More explicitly, let \(\kappa_1,\ldots,\kappa_{n_p}\) be this initial
list of momenta, and let \(\iota:\mathcal V_{\ell'}\to\mathcal V_{n_p}\)
insert the normalized double singlets in its pair slots.
If \(\mathfrak S_{\ell'm'}\) consists of the permutations \(\pi\)
of labels into slots satisfying
\(\pi(A_a)<\pi(B_a)\) for every labeled low--high pair \(A_aB_a\),
then
\[
 E_X^{\ell'm'}(k,u)
   =\sum_{\pi\in\mathfrak S_{\ell'm'}}
       e^{\,i\sum_{j=1}^{n_p}\kappa_jX_{\pi(j)}}A_\pi(k,u)\iota,
\]
where \(A_\pi\) is the product of adjacent interchange matrices
along a reduced sequence producing \(\pi\), and \(A_{\rm id}=I\).
Coefficients reversing a pair are zero.  The inverse and braid
relations make this prescription independent of the chosen sequence
of adjacent interchanges.  At a removable singularity we use its
continuous value.  The initial order of the whole blocks can be chosen
arbitrarily; the block-interchange formulas proved below describe the
change of this convention.

\begin{lemma}[Line spectral resolution]
\label{lem:bethe-line}
For every fixed integer \(n_p=2M+\ell\geq0\), every \(c>0\), and every
\(s>0\), the heat kernel of \(H_c\) on the ordered line chamber is
\begin{equation}
\begin{split}
 e^{-sH_c}(X,Y)
  ={}&\sum_{\ell'+2m'=n_p}
       \frac{(2c)^{m'}}{\ell'!\,m'!}
       \int_{\mathbb R^{\ell'+m'}}
       E_X^{\ell'm'}(k,u)E_Y^{\ell'm'}(k,u)^\dagger\\
    &\quad{}\times
       \exp\left\{-s\left(
           \sum_{a=1}^{\ell'}k_a^2+
           2\sum_{a=1}^{m'}u_a^2-\frac{m'c^2}{2}
                    \right)\right\}
       \prod_{a=1}^{\ell'}\frac{dk_a}{2\pi}
       \prod_{a=1}^{m'}\frac{du_a}{2\pi}.
\end{split}
\label{eq:B10}
\end{equation}
Empty products and the vacuum term have their usual value \(1\).
The transform defined by the kernels in \eqref{eq:B10}, with the
square roots of the displayed weights, is an isometry.  Moreover,
\(H_c+Mc^2/2\geq0\), and its quadratic form is the integral of
\[
  \sum_{a=1}^{\ell'}k_a^2+
  2\sum_{a=1}^{m'}u_a^2+\frac{(M-m')c^2}{2}
\]
against the squared modulus of this transform.
\end{lemma}

\begin{proof}
The proof has three parts.  The scattering identities show that taking
a pair residue creates no further bound cluster.  A heat-kernel formula
on separated contours then supplies the initial condition.  Finally,
lowering those contours produces \eqref{eq:B10}, with its normalization
determined by the particle occupying the first slot at each step.

\emph{Scattering of pairs and cancellation of higher residues.}
Let \(\Pi_{\varepsilon\eta}\), with
\(\varepsilon,\eta\in\{+1,-1\}\), be the joint spectral projections of
\(P^L,P^R\) on two slots.  The interchange matrix has the decomposition
\[
 \check R(x)=
 \frac{x-ic}{x+ic}\Pi_{++}
 +\frac{x+ic}{x-ic}\Pi_{--}
 -\Pi_{+-}-\Pi_{-+}.
\]
In particular,
\[
 \mathop{\rm Res}_{x=ic}\check R(x)=2ic\,\Pi_{--},
 \qquad \check R(0)=-I.
\]
Here \(\Pi_{--}\) is precisely the double-singlet projection.  Taking
this residue produces a low--high pair in the order specified above.
Multiplication of two interchanges transports a pair past a single
and multiplies its color vector by
\begin{equation}
 S_{21}(x)=
 \frac{(x+ic/2)(x-3ic/2)}
      {(x-ic/2)(x+3ic/2)}.
\label{eq:bethe-pair-single}
\end{equation}
In this formula \(x\) is the center of the left block minus the
momentum of the right block.  If the single is on the left, the same
formula holds with \(x\) equal to its momentum minus the pair center.
Moving a pair past another pair gives
\begin{equation}
                       S_{22}(x)=\frac{x-2ic}{x+2ic},
\label{eq:bethe-pair-pair}
\end{equation}
where again \(x\) is the left center minus the right center.
Both scalars have modulus one for real \(x\).

Here are explicit checks, including the partial exchanges that will
occur during contour motion.  In one binary index, for a pair \(AB\)
followed by a single \(C\), set
\[
 E(q)=s_{12}\otimes q_3,\qquad
 H(q)=q_1\otimes s_{23},\qquad
 F(q)=s_{13}\otimes q_2=E(q)+H(q).
\]
For an operator or vector \(V\) in one binary index, use
\(V^{[2]}=V^L\otimes V^R\).  Put
\[
 D(x)=(x-ic/2)(x+3ic/2),\qquad
 T(x)=(x-ic/2)E+(x+ic/2)H.
\]
With \(x=u-k\), the coefficients of the three shuffles \(ABC,ACB,CAB\)
are, respectively,
\[
            E^{[2]},\qquad
            \frac{T(x)^{[2]}}{D(x)},\qquad
            S_{21}(x)H^{[2]}.
\]
These identities follow by applying the two permutation factors in
\(\check R\) to the singlet.  For a single initially on the left,
interchange \(E\) and \(H\), and take \(x=k-u\).  Thus neither partial
nor complete exchanges have a pole at \(x=3ic/2\).  At \(x=ic/2\),
the last two residues in the pair-first order are
\((ic/2)H^{[2]}\) and \(-(ic/2)H^{[2]}\).
The elementary momentum then equals the lower member of the pair,
so their plane phases agree and the residues cancel.  In the
single-first order the corresponding equality is with the upper
member.  This is also the cancellation expressed by
\(\check R(0)=-I\).

For two pairs \(AB,CD\), both in low--high order, put \(x=u-v\).
In one binary index set
\[
        e=s_{12}s_{34},\qquad f=s_{14}s_{23},\qquad
        h=s_{13}s_{24}=e+f.
\]
The coefficients of the six allowed shuffles are
\[
\begin{array}{c|c}
 ABCD & e^{[2]}\\[2pt]
 ACBD & (xh+icf)^{[2]}/\{x(x+2ic)\}\\[2pt]
 ACDB & (x^2+c^2)f^{[2]}/\{x(x+2ic)\}\\[2pt]
 CABD & (x^2+c^2)f^{[2]}/\{x(x+2ic)\}\\[2pt]
 CADB & (xh-icf)^{[2]}/\{x(x+2ic)\}\\[2pt]
 CDAB & (x-2ic)e^{[2]}/(x+2ic).
\end{array}
\]
This table can equally be obtained from \(P_{12}=-1\) on a singlet
and
\(1-P_{12}-P_{13}-P_{23}+P_{12}P_{23}+P_{23}P_{12}=0\)
on three binary slots.  It proves \eqref{eq:bethe-pair-pair} and
shows that the apparent poles at \(x=ic\) and \(x=2ic\) are absent.
At \(x=0\), the four middle residues are, in their listed order,
\[
       \frac{ic}{2}f^{[2]},\quad
       -\frac{ic}{2}f^{[2]},\quad
       -\frac{ic}{2}f^{[2]},\quad
       \frac{ic}{2}f^{[2]}.
\]
The phases agree in the pairs \(ACBD,ACDB\) and \(CABD,CADB\),
because \(B\) and \(D\) have the same momentum.  The cancellation
therefore holds both in the whole wave and in either part in which
the first slot is required to contain the low member of a specified
pair: the part starting with \(A\) consists of the first three
shuffles, and the part starting with \(C\) of the last three.

\emph{The heat kernel on separated contours.}
Let \(E_X^{n_p,0}(k)\) be the elementary-wave matrix, and choose
horizontal contours
\(\Gamma_j=\mathbb R+i\alpha_j\), where
\[
 \alpha_{n_p}=0,\qquad \alpha_j-\alpha_{j+1}>c.
\]
Consider
\begin{equation}
 \mathcal K_s(X,Y)=
 \int_{\Gamma_1\times\cdots\times\Gamma_{n_p}}
 E_X^{n_p,0}(k)
 e^{-i\sum_j k_jY_j-s\sum_j k_j^2}
 \prod_j\frac{dk_j}{2\pi}.
\label{eq:bethe-separated-contours}
\end{equation}
Gaussian decay makes this integral and its spatial derivatives
convergent on compact sets.  Each plane wave solves the free heat
equation in the chamber, and the interchange relation imposes the
contact conditions on each face.

To verify the initial condition, one needs the support of the
Fourier transforms of the rational coefficients.  For
\(\operatorname{Im}x>c\), the preceding projector decomposition gives
the concrete formula
\[
 \check R(x)=P^LP^R+
 2c\int_0^\infty e^{ixt}
       \bigl(e^{ct}\Pi_{--}-e^{-ct}\Pi_{++}\bigr)\,dt.
\]
Accordingly, an inversion of labels \(i<j\) contributes a
nonnegative variable \(t_{ij}\), possibly an atom at zero, and a phase
\(e^{i(k_i-k_j)t_{ij}}\).  If a permutation puts label \(i\) at
position \(p(i)\), its zero-time contribution can be supported only at
\begin{equation}
 Y_i=X_{p(i)}+
       \sum_{\substack{j>i\\p(j)<p(i)}}t_{ij}
       -\sum_{\substack{h<i\\p(h)>p(i)}}t_{hi}.
\label{eq:bethe-inversion-support}
\end{equation}
For a nonidentity permutation, let \(r\) be the first nonfixed
position, and let \(j>r\) be the label occupying it.
Label \(r\) occurs at some position \(p(r)>r\).
No smaller label contributes a negative shift to label \(r\);
no larger label contributes a positive shift to label \(j\).
Thus \eqref{eq:bethe-inversion-support} would imply
\[
                   Y_r\geq X_{p(r)}>X_r\geq Y_j,
\]
which is incompatible with \(Y\in\mathcal C_{n_p}\).
Only the identity permutation remains, and its coefficient is the
identity color operator.  It supplies the delta initial condition.

The support argument also supplies the estimates needed to justify
this limiting calculation and heat-kernel uniqueness.  Let \(d\)
be the shift vector on the right side of
\eqref{eq:bethe-inversion-support}, so that
\(z_i=X_{p(i)}+d_i\).  The shifts belong to the cone generated by
\(e_i-e_j\), \(i<j\).  This cone is pointed, and more explicitly
\[
 \sum_{i<j}t_{ij}
 \leq \sum_{i<j}(j-i)t_{ij}
 =-\sum_i i\,d_i
 \leq C_{n_p}|d|.
\]
Furthermore,
\[
 \sum_i X_{p(i)}d_i
   =\sum_{\substack{i<j\\p(j)<p(i)}}
        t_{ij}(X_{p(i)}-X_{p(j)})\geq0,
 \qquad |z|^2\geq |X|^2+|d|^2.
\]
The measures in the transform formula have at most exponential
growth in \(\sum t_{ij}\).  Gaussian convolution therefore dominates
them and gives spatial decay for compactly supported initial data.
It also justifies convergence to the stated initial condition,
first for data supported strictly inside the chamber and then by
density.  The contact form is lower bounded by the hyperplane trace
inequality, so uniqueness of its \(L^2\) heat evolution identifies
\(\mathcal K_s\) with \(e^{-sH_c}\).

\emph{Lowering the contours and reconstructing the adjoint wave.}
Process the momenta from last to first.  The induction assertion can be
stated explicitly.  After the last \(n_p-h\) slots have been processed,
sum over \(\ell'+2m'=n_p-h\), with weight
\((2c)^{m'}/(\ell'!m'!)\).  The first \(h\) elementary momenta
\(\xi_1,\ldots,\xi_h\) remain on their original high contours; the
suffix momenta \(k\) and pair centers \(u\) are real.  The wave factor
in the integrand is
\[
 E_X^{h+\ell',m'}(\xi,k,u)
 \left[I_{\mathcal V_h}\otimes
       E_{(Y_{h+1},\ldots,Y_{n_p})}^{\ell'm'}(k,u)^\dagger\right]
 e^{-i\sum_{a=1}^h\xi_aY_a}.
\]
Its Gaussian energy is
\(\sum\xi_a^2+\sum k_a^2+2\sum u_a^2-m'c^2/2\).
Here the \(X\)-wave uses the first \(h\) elementary slots followed by
the suffix blocks as its initial order; its finite meromorphic formula
also defines it for the complex variables \(\xi\).  Thus \(h=n_p\)
is \eqref{eq:bethe-separated-contours}, whereas \(h=0\) is precisely
\eqref{eq:B10}.

To pass from \(h\) to \(h-1\), decompose the new suffix wave according
to its first occupied slot.  For each elementary carrier or labeled pair
\(a\), let \(E_Y^{[a]}\) be the partial sum in which that slot contains,
respectively, the elementary momentum or the pair's low member.  These
possibilities are disjoint and exhaust the allowed shuffles, so
\(E_Y=\sum_a E_Y^{[a]}\).  The ordinary contour integral will supply
the elementary-first partial adjoints, and its residues will supply the
pair-first partial adjoints.
When block parameters \(z\) leave the real contours, an adjoint wave is
continued as the meromorphic function \(E_Y(\overline z)^\dagger\),
which agrees with \(E_Y(z)^\dagger\) for real \(z\).  The same convention
applies to each partial adjoint.

Let \(k=\xi_h\) be the next momentum.  The integral left after lowering
its contour to the real line provides the adjoint-wave terms whose
first suffix slot is an elementary momentum.  To see this with a
fixed elementary carrier, put that carrier first among the suffix
blocks.  The unitary elementary and whole-block interchanges move
the corresponding coefficient from the \(X\)-wave to the
\(Y\)-wave.  Symmetrizing the elementary variables then supplies
all elementary choices for the first slot.

The residues occur at \(k=p+ic\), where \(p\) is an elementary
momentum in the old suffix.  Put \(p\) first in that suffix before
computing the residue.  A downward contour shift contributes
\(-i\) times the residue, so the residue \(2ic\,\Pi_{--}\) contributes
the factor \(2c\) and inserts a normalized double singlet.
The momenta at \(X\) are \(p,k\), in low--high order, with center
\(u=p+ic/2\).  Lower \(u\) from this contour to the real line.
On the \(Y\) side the new first slot uses the adjointed phase of
the low member, namely the phase with the high momentum \(k\).
The other member is the old suffix variable \(p\).  The old
suffix exchanges are continued with this \(p\), and hence give
exactly the part of \(E_Y^\dagger\) with that pair's low member
first.

For completeness, the pole checks justifying both contour motions
are as follows.  When \(k\) descends past an existing pair, the
three-slot formulas have no pole at center difference \(3ic/2\);
the possible pole at \(ic/2\) cancels between the partial and
complete exchanges whose phases agree.  There is therefore no
three-particle residue.  When the newly formed center \(u\) is
lowered, the same cancellation applies to the \(X\)-wave.  In the
restricted adjoint \(Y\)-expression the low member of the new pair
has already been put first.  The crossing that remains uses only
\(p=u-ic/2\).  Its imaginary part moves from \(0\) to \(-c/2\),
and no elementary interchange pole at difference \(\pm ic\) is
crossed.  This includes an elementary momentum on the old real
contour, by continuation with an indentation if necessary.
Explicitly, the adjoint of the pair-low-first three-slot sum has
denominator \((x+ic/2)(x-3ic/2)\), with \(x=u-k_{\rm single}\),
which has no zero in \(0\leq\operatorname{Im}x\leq c/2\).
For two pairs, the four-slot table proves regularity at coincident
centers even in the restricted first-slot expression.
The only other possible center pole is at \(-2ic\) for the
\(X\)-coefficients and at \(2ic\) for the continued adjoint
coefficients, both outside the strips being traversed.

These checks remain valid with other variables inserted.  At a
generic point of a candidate pole hyperplane, use the braid and
whole-block identities to bring the two relevant blocks together.
For fixed occupied slots, inserting the other momenta multiplies
the two cancelling residues by the same nonsingular interchange
operators.  If the first slot is restricted, put its block first
before this operation.  The earlier, unprocessed contours stay
separated from all these motions.  Finally, write the finite sums
over a common rational denominator.  Removal of the poles at
generic points of every denominator hyperplane removes the
corresponding factors, so the removals extend across their
intersections.  Distinct real parts or indented contours may
therefore be used for the residue computations and subsequently
removed.  On bounded \(X,Y\), small contour circles about removable
singularities give uniform bounds for the finite wave sums.
Together with the Gaussian, these bounds justify the integrations
and deformations just made.

It remains to give every partial adjoint the same weight.  Changing
the initial block order multiplies both waves by the same unitary
block-interchange operator, so the product
\(E_X(E_Y^{[a]})^\dagger\) is independent of that choice.  We may
therefore put each first carrier in front before symmetrizing the
labels.  For a final sector \((\ell',m')\), the
elementary-first contribution comes from
\((\ell'-1,m')\) in the old suffix and has weight
\[
 \frac{(2c)^{m'}}{(\ell'-1)!\,m'!}
       =\ell'\frac{(2c)^{m'}}{\ell'!\,m'!}.
\]
It accounts for the \(\ell'\) choices of elementary first carrier.
The pair-first contribution comes from
\((\ell'+1,m'-1)\).  There are \(\ell'+1\) elementary choices
for \(p\), and the residue supplies \(2c\); its weight is therefore
\[
 2c(\ell'+1)
 \frac{(2c)^{m'-1}}{(\ell'+1)!\,(m'-1)!}
       =m'\frac{(2c)^{m'}}{\ell'!\,m'!}.
\]
It accounts for the \(m'\) choices of pair first carrier.
After symmetrization, each of the \(\ell'+m'\) partial adjoints has
weight \((2c)^{m'}/(\ell'!m'!)\).  Their sum is the full adjoint wave.
This closes the induction and proves \eqref{eq:B10}.

Define the transform explicitly by
\[
 (\mathcal F_c f)_{\ell'm'}(k,u)=
 \left(\frac{(2c)^{m'}}{\ell'!\,m'!}\right)^{1/2}
 \int_{\mathcal C_{n_p}}E_Y^{\ell'm'}(k,u)^\dagger f(Y)\,dY,
\]
initially for compactly supported \(f\).  Its target is the direct
sum of the \(L^2\) spaces with the product measures \(dk\,du/(2\pi)\)
in \eqref{eq:B10} and color space \(\mathcal V_{\ell'}\).
Writing
\[
 \lambda_{\ell'm'}(k,u)=
 \sum k_a^2+2\sum u_a^2+\frac{(M-m')c^2}{2}\geq0,
\]
the quadratic identity from \eqref{eq:B10} is
\[
 \langle f,e^{-s(H_c+Mc^2/2)}f\rangle
 =\sum_{\ell'+2m'=n_p}\int
       e^{-s\lambda_{\ell'm'}}|(\mathcal F_cf)_{\ell'm'}|^2.
\]
As \(s\downarrow0\), strong continuity on the left and monotone
convergence on the right give \(\|\mathcal F_cf\|=\|f\|\).
The transform thus extends isometrically to the full Hilbert
space.  The same identity then implies
\(\|e^{-s(H_c+Mc^2/2)}\|\leq1\), and hence
\(H_c+Mc^2/2\geq0\).  Taking the increasing limit of
\((1-e^{-s\lambda})/s\) gives, by the spectral theorem,
the stated quadratic-form identity and its domain.
The reconstruction \(\mathcal F_c^*\mathcal F_c=I\) is all that
will be needed; surjectivity onto the direct sum is unnecessary.
\end{proof}

\subsection{Pair concentration at strong coupling}
\label{subsec:bethe-strong-coupling}

The spectral resolution contains only single particles and singlet pairs;
the residue cancellations have excluded additional binding channels.
Its form identity gives a penalty \(c^2/2\) for every missing pair.
To use this information on a circle, we also need control of the spatial
shape of the pair waves, uniformly in their real center momenta.

For a shuffle contributing to \(E_X^{\ell'm'}\), draw the closed
interval between the two slots occupied by each pair.  Combine
overlapping intervals into connected components, and include every
single whose slot lies inside such a component.  A single outside
all pair intervals forms a component of one slot.  Group the
shuffles for which every component has the same set of slots and
the same set of carriers (elementary momenta and labeled pairs).
Let \(E_{X,\mathfrak g}^{\ell'm'}\) denote the sum for one such group.

\begin{lemma}[Decay across pair spans]
\label{lem:bethe-shuffle-decay}
At fixed \(n_p\), there are constants \(C_{n_p}<\infty\) and
\(a_{n_p}>0\), independent of \(c>0\) and the real momenta, such that
every group just defined satisfies
\begin{equation}
 \bigl\|E_{X,\mathfrak g}^{\ell'm'}(k,u)\bigr\|
 \leq C_{n_p}
 \exp\left\{-a_{n_p}c
       \sum_{\text{components }\mathcal I}
       \bigl(X_{\max\mathcal I}-X_{\min\mathcal I}\bigr)\right\}.
\label{eq:B11}
\end{equation}
The norm is the operator norm after summing the shuffles in the
group.
\end{lemma}

\begin{proof}
Divide all momenta by \(c\).  A shuffle's plane phase has modulus
\[
 \exp\left\{-\frac c2
       \sum_{\text{pairs }a}
           (X_{\text{high slot of }a}
                 -X_{\text{low slot of }a})\right\}.
\]
The sum of the pair spans is at least the sum of the component
spans.  Rational interchange coefficients are uniformly bounded
away from their pole hyperplanes.  Near real centers, their only
apparent singularities are coincident pair centers.  The
four-slot cancellations in the preceding proof are internal to
the present groups: switching the equal high members preserves
the slots and carriers of each overlapping component.
More explicitly, the cancelling shuffles are
$ACBD,ACDB$ and $CABD,CADB$.  Within either pair of shuffles,
the union of the spans belonging to $AB$ and $CD$ is the same
interval, from their earliest low member to their latest high member.
Inserting further slots does not change that equality.  The connected
components obtained after adjoining all other pair spans, and their
sets of carriers and included singles, are therefore unchanged.
Thus every group sum is holomorphic through these real
hyperplanes.

Here is a uniform Cauchy estimate for this holomorphic sum.
Choose a sufficiently small number depending only on \(n_p\).
Among finitely many successively much smaller scales, one can
choose a scale with no distance between real divided momenta in
the intervening annulus: there are only finitely many pairwise
distances.  Cluster the arguments at the smaller scale.
Arguments in different clusters are then separated by much more
than the chosen scale.  Within each cluster, put the variables
on small circles about their original real values,
using distinct radii separated by fixed multiples of that scale.
All radii can be kept below the initially chosen small number.
On the product of these circles, equal-center denominators are
bounded away from zero by a positive \(n_p\)-dependent constant;
denominators with nonzero imaginary offsets are also separated
from zero.  Every rational coefficient is consequently bounded
there by a constant depending only on \(n_p\).  Cauchy's formula
recovers its holomorphic group sum at the original real arguments.

To control the plane phases on those circles, factor from each
component the common phase evaluated at its leftmost position.
The carrier set of that component is fixed throughout the group,
so the factored phase is the same in every term and has modulus
one at real momenta.  Perturbing the remaining divided momenta by
the small circle radii increases their modulus by at most
\(\exp\{\varepsilon C_{n_p}c\sum_{\mathcal I}
\operatorname{span}(\mathcal I)\}\).
Choose the initial upper bound on the radii small enough that
this uses less than half the original pair-gap decay.
The resulting Cauchy estimate proves \eqref{eq:B11}.
\end{proof}

We next specify the compactness statement on a circle.  Fix its
length \(W_0>0\).  Choose finitely many seams and smooth nonnegative,
permutation-invariant cutoffs \(\vartheta_\alpha\), with
\(\sum_\alpha\vartheta_\alpha^2=1\), such that all positions in the
support of \(\vartheta_\alpha\) stay a positive distance from the
corresponding seam.  Such a choice is possible because, at fixed
particle number, among sufficiently many candidate seams at least
one stays a fixed distance from every position.  Cutting at that
seam identifies each localized function with a compactly supported
function on a line chamber.

In each chamber let \(\mathfrak p\) be a partition of its ordered
slots into \(M\) adjacent pairs and, when \(\ell=1\), one single.
For a pair occupying slots \(a,a+1\), use coordinates
\[
 q=\frac{X_a+X_{a+1}}2,\qquad
 z=c(X_{a+1}-X_a)>0.
\]
Keep the single coordinate, if any, unchanged.
The coordinate Jacobian is \(c^{-M}\), so the corresponding
isometric change of \(L^2\) coordinates multiplies a wavefunction
by \(c^{-M/2}\).  Pair centers and the single coordinate are
called separated when all their mutual distances are positive.
The sets of nonseparated center coordinates have measure zero.
For \(\delta>0\) and \(Z<\infty\), a retained pair chart means the
part of these coordinates on which all center distances exceed
\(\delta\) and \(0<z_a<Z\) for every pair.
The corresponding regions for distinct pair assignments are
disjoint once \(c\) is sufficiently large in terms of \(\delta,Z\).
Convergence in these charts will always be followed by their
joint exhaustion as \(\delta\downarrow0\) and \(Z\uparrow\infty\).
It does not mean that the whole wavefunction has an \(L^2\) limit
under every possible pairing coordinate map.

\begin{lemma}[Compactness in pair coordinates]
\label{lem:bethe-pair-compactness}
Fix \(W_0>0\), \(n_p=2M+\ell\), and \(C<\infty\).
Let \(c_j\to\infty\) and let \(\psi_j\) be normalized vectors in
the form domain of the circle Hamiltonian such that
\[
 \left\langle\psi_j,
        (H_{c_j}+Mc_j^2/2)\psi_j\right\rangle\leq C.
\]
After passage to a subsequence, the localized functions in the
pair coordinates just defined converge in \(L^2\) on every fixed
retained chart.  For each seam and pair assignment these limits
are compatible under enlargement of the chart and have the form
\begin{equation}
 f_{\alpha,\mathfrak p}(q,v)\,
 \exp\left\{-\frac12\sum_{a=1}^M z_a\right\}
 \bigotimes_{a=1}^M
        (s_a^L\otimes s_a^R).
\label{eq:B12}
\end{equation}
Here \(v\) and its color are present only if \(\ell=1\); the
coefficient \(f_{\alpha,\mathfrak p}\) includes that color.
As \(\delta\downarrow0\) and \(Z\uparrow\infty\), the sum of the
squared norms of these chart limits tends to \(1\), counted with
the fixed square partition of unity.
There is no norm lost away from the pair coincidence manifolds
or at their intersections.  Inner products of two sequences
converging in this sense converge to the corresponding sums of
inner products of their limits.
\end{lemma}

\begin{proof}
The form localization identity gives
\[
 \sum_\alpha
 \left\langle\vartheta_\alpha\psi_j,
      (H_{c_j}+Mc_j^2/2)\vartheta_\alpha\psi_j\right\rangle
 =
 \left\langle\psi_j,
      (H_{c_j}+Mc_j^2/2)\psi_j\right\rangle
 +\sum_\alpha\|\,|\nabla\vartheta_\alpha|\psi_j\|^2.
\]
The error is bounded independently of \(j\), and every localized
term is nonnegative by Lemma~\ref{lem:bethe-line}.
It suffices, therefore, to prove the assertion for one localized
function on the line, with bounded norm and bounded shifted form
energy.  Fix a bounded chamber box containing the supports of
all these localized functions.  In what follows, restrict the
output of every reconstruction kernel to this box.  This is an
\(L^2\) contraction and changes none of the localized functions;
it is needed because a momentum-truncated reconstruction need
not itself have compact support.

Apply the isometric transform of Lemma~\ref{lem:bethe-line}.
Restrict all elementary momenta and pair centers to a fixed
bounded set, say absolute value at most \(R\).
The form identity bounds the squared norm discarded by this
restriction by \(C'/R^2\).  Its reconstruction has at most this
norm because \(\mathcal F_c^*\) is a contraction.
Likewise, all sectors with fewer than \(M\) pairs have
\(\lambda_{\ell'm'}\geq c^2/2\); their total squared transform
norm is at most \(2C'/c^2\).  Thus only the sector with \(M\)
pairs and \(\ell\) singles can survive, and its momenta may first
be kept in a fixed bounded set.

For that sector the reconstruction kernel contains the factor
\((2c)^{M/2}/\sqrt{\ell!\,M!}\).  Consider a group with overlapping
pair spans or with a single inside a pair span.
If a component contains \(p\) pairs and \(s_0\) singles, it has
\(2p+s_0\) slots and \(2p+s_0-1\) internal gaps.
For a component consisting of one adjacent pair this is one gap.
For an overlapping component or one containing an inserted
single it is at least \(p+1\) gaps.
Squaring \eqref{eq:B11} and integrating the ordered coordinates
therefore gives one factor \(O(c^{-1})\) for every such gap.
Over all components there are at least \(M+1\) decaying gaps in
these groups.  The squared kernel normalization is \(O(c^M)\),
so their squared Hilbert--Schmidt norms on the momentum cutoff
are \(O(c^{-1})\).  The remaining component positions range over
a fixed bounded set.  The bounded momentum volume and the
Cauchy--Schwarz inequality now show that all these groups have
vanishing reconstructed \(L^2\) norm, uniformly for bounded
transform coefficients.

Every surviving shuffle is a permutation of whole adjacent
blocks.  In coordinates for its adjacent-pair assignment, its
normal plane factor is
exactly \(\exp(-\sum z_a/2)\) and its color factor is the product
of double singlets in \eqref{eq:B12}.  The coefficients between
whole blocks are the explicit scalars
\eqref{eq:bethe-pair-single}--\eqref{eq:bethe-pair-pair}, whose
limits on a fixed momentum cutoff exist as \(c\to\infty\).
The rescaled Jacobian cancels the factor \(c^{M/2}\) in the
kernel normalization.  The fixed center-coordinate volume,
the finite momentum cutoff, and the exponential normal factor
give an integrable dominating square.
The reconstruction kernels for this assignment consequently
converge in Hilbert--Schmidt norm on every retained chart.
On a chart belonging to another assignment, their contribution
tends to zero: bounded rescaled gaps for both assignments would
force two distinct block centers to approach each other, whereas
the retained chart separates those centers by \(\delta\).
The exponential gap bound controls the remaining contribution.

This kernel convergence implies compactness even though the
transform coefficients depend on \(j\).  On the fixed cutoff
their \(L^2\) norms are bounded; extract a weakly convergent
subsequence.  The limiting Hilbert--Schmidt operator is compact,
and convergence to it in operator norm then gives strong
convergence of the reconstructed functions.
The limit has the asserted normal and color factors.

We verify that this argument retains the whole norm.
Restrict first to centers separated by at least \(\delta>0\)
and to \(0<z_a<Z\).  Different adjacent-pair assignments then
occupy disjoint configuration regions for all sufficiently
large \(c\).  The retained region is the union over those
assignments; each whole-block kernel is estimated in the
coordinates of its own assignment.
For the bounded-momentum kernels, the squared
norm excluded by center distances below \(\delta\) is uniformly
small as \(\delta\downarrow0\): the kernels are uniformly
bounded in the center coordinates, and these collision
neighborhoods have vanishing volume.
The norm excluded by \(z_a>Z\) tends uniformly to zero as
\(Z\to\infty\), by the exponential normal bound.
The groups already discarded cover the overlaps and inserted
singles, and have vanishing norm.
Thus neither intersections of pair coincidence manifolds nor
their complement retain additional norm.

Finally remove the momentum cutoff using the uniform
\(C'/R^2\) bound, and take a diagonal subsequence over the
finitely many seam charts and pair assignments and a countable
exhaustion of the separation and gap cutoffs.
The exact identity \(\sum_\alpha\vartheta_\alpha^2=1\)
then gives the asserted total norm.
Applying the same cutoffs to two sequences, strong convergence
on the retained charts gives convergence of their inner
products; Cauchy--Schwarz controls all discarded parts.
\end{proof}

The needed ground-state test is now an overlap argument.  The center
wave need not be characterized by a separate effective Hamiltonian:
strict positivity of its limiting shape suffices.

\begin{corollary}[Ground-state test by positive overlap]
\label{cor:bethe-ground-test}
Fix a circle length \(W_0>0\) and a particle number
\(n_p=2M+\ell\), with \(\ell\in\{0,1\}\), and work in the
fixed-weight sector balanced in each binary index.  Suppose a normalized eigenbranch
has energy \(-Mc^2/2+O(1)\) as \(c\to\infty\).  Suppose also that,
after one common phase choice for all charts, every subsequential limit
in the pair charts has the form \eqref{eq:B12}.  Require its scalar
center coefficient to be strictly positive almost everywhere that the
limiting chart cutoff is positive, in the static-sign convention of
Lemma~\ref{lem:brownian-ground}; when \(\ell=1\), the remaining single
has the majority color in each index.
Then the branch is the ground state for all sufficiently large \(c\).
\end{corollary}

\begin{proof}
Trial states consisting of separated, exponentially bound
pairs and smooth center wavefunctions give a circle ground-state
energy at most \(-Mc^2/2+O(1)\).  One can cut off the pair gaps
inside fixed disjoint intervals: the cutoff error is
exponentially small in \(c\), and the center kinetic energy is
bounded.  The true normalized ground state therefore satisfies
Lemma~\ref{lem:bethe-pair-compactness}.
It is the simple ground state in this sector, is nonnegative
in the static-sign convention, and its limits
retain total norm \(1\).  Their normal factors are
\eqref{eq:B12}; in a sector balanced in each binary index, the
remaining single, if present, has the majority color in each
index.  A nonzero nonnegative limit consequently has strictly positive
inner product with the strictly positive limiting center coefficient of
the eigenbranch.

If the two eigenstates were different for an unbounded sequence
of couplings, self-adjointness and the simplicity of the ground
state would make their inner products zero on that sequence.
Pass to a simultaneous compactness subsequence.  The inner-product
convergence just proved gives a strictly positive limit, a
contradiction.  This identifies the strong-coupling branch from
its wavefunction and energy, without assuming a string
description of the ground state on a finite circle.
\end{proof}


\subsection{Center equations and uniform root estimates}

We return to $c=2$.  A pair will have momenta close to $u_j-i,u_j+i$ and
auxiliary variable close to $u_j$; when $\ell=1$ there is also a single
momentum $v$.  Before solving for the exponentially small changes in the
pair spacings, we solve the real equations for the pair centers.  The
construction runs along the entire ray $\mathcal W\geq W$ with $n_p$ fixed.
Its large-$\mathcal W$ end corresponds, after rescaling the circle, to
the strong-coupling limit of Corollary~\ref{cor:bethe-ground-test}.
The estimates must remain uniform down to $\mathcal W=W$.

At exact pair spacings, multiply the two momentum equations in
\eqref{eq:B8} and eliminate their internal auxiliary factors by the
auxiliary equation \emph{before} passing to the limit.  With consecutive
pair quantum numbers centered at zero, this gives
\begin{equation}\label{eq:B13}
 \begin{split}
 2\mathcal W u_j
   &=\sum_{t\ne j}\left[
       \pi\operatorname{sgn}(j-t)
       -2\arctan\frac{u_j-u_t}{4}\right]
       +\ell\varphi(u_j-v),\\
 \mathcal W v&=\sum_{j=1}^M\varphi(v-u_j),
       \qquad
 \varphi(x)=2\left(\arctan x-\arctan\frac{x}{3}\right).
 \end{split}
\end{equation}
The second equation is present only when $\ell=1$.  In particular these
equations are obtained from the full Bethe equations, not by omitting a
vanishing internal factor.

\begin{lemma}[Counting density and the single-particle Jacobian]
\label{lem:bethe-centers}
Under the parameter assumptions at the start of this Section, for every
$\mathcal W\geq W$, \eqref{eq:B13} has a reflection-symmetric solution
with $v=0$ and $u_1<\cdots<u_M$.  With $v$ fixed, the pair solution is
unique.  Define
\[
 Z_c(x)=\frac{x}{\pi}
    +\frac{1}{\pi\mathcal W}\sum_{j=1}^M
           \arctan\frac{x-u_j}{4}
    -\frac{\ell}{2\pi\mathcal W}\varphi(x),
 \qquad d(x)=Z_c'(x).
\]
Then $Z_c$ is odd and
\[
 d(x)=\frac1\pi+\frac1{\mathcal W}\sum_jK(x-u_j)
             -\frac\ell{\mathcal W}(K_1-K_3)(x),
 \qquad
 Z_c(u_j)=\frac{j-(M+1)/2}{\mathcal W}.
\]
Let $B_c>0$ be determined by $Z_c(B_c)=M/(2\mathcal W)$.
The following bounds hold uniformly:
\begin{enumerate}[label=\textup{(\roman*)}]
\item $d$ is bounded below by a positive constant;
      $B_c=O((1+b)^2)$; $d$ and any fixed number of its derivatives
      are bounded by $\operatorname{poly}(b)$.  Neighboring centers have
      separation at least $1/(\mathcal W\operatorname{poly}(b))$.
\item On $(-B_c,B_c)$,
\begin{equation}\label{eq:B14}
 d=\rho_{B_c}
  -\frac\ell{\mathcal W}(1-K*)_{B_c}^{-1}(K_1-K_3)
  +O\!\left(\frac{\operatorname{poly}(b)}{\mathcal W^2}\right).
\end{equation}
The inverse function in the middle term has bounded supremum and
$L^1$ norm, independently of $B_c$.
\item $B_c\leq B_*+O(1)$.  At $\mathcal W=W$ one also has
      $B_c\geq B_*-O(1)$.
\item The full Jacobian of \eqref{eq:B13}, including the single equation
      when present, is invertible.  Its inverse has row sums bounded
      by a fixed power of $\mathcal W$.
\end{enumerate}
\end{lemma}

\begin{proof}
Move the increasing pair terms in the first equation to the left.  For
fixed $v$, these equations are the gradient of a strictly convex function:
the quadratic part has Hessian $2\mathcal W I$, the pair interactions
give a nonnegative weighted graph Laplacian, and the term involving
$\varphi$ has uniformly bounded diagonal curvature.  The Hessian is
therefore strictly positive.  Coercivity gives existence and uniqueness.
The increasing right-hand sides $\pi(2j-M-1)$ imply the ordering of the
$u_j$; this follows also by interchanging two variables in the convex
minimization problem.  Reflection reverses the index order and changes
every sign, so uniqueness gives $u_j=-u_{M+1-j}$ when $v=0$.
Oddness of $\varphi$ then proves the single equation.

The displayed formulas for $Z_c$ and $d$ follow by division and
differentiation of \eqref{eq:B13}.  Since $K\geq0$, its only possibly
negative term is $O(\mathcal W^{-1})$, giving the lower bound for $d$.
The total quantile interval has length $M/\mathcal W=O(b^2)$, which gives
the first bound for $B_c$.  The derivatives of each Cauchy kernel are
bounded, and $M/\mathcal W=O(b^2)$, proving the derivative bounds.  The
quantile increments are $1/\mathcal W$, so the upper bound for $d$ gives
the stated separation.

For a smooth function of $u$, summation over the centers is midpoint
summation after the change of variable $s=Z_c(u)$.  The inverse change of
variable and its fixed derivatives have polynomial bounds because $d$
is bounded below.  Applying this rule to $K(x-u)$, uniformly in $x$, gives
\[
 \frac1{\mathcal W}\sum_jK(x-u_j)
   =\int_{-B_c}^{B_c}K(x-y)d(y)\,dy
       +O\!\left(\frac{\operatorname{poly}(b)}{\mathcal W^2}\right).
\]
The inverse of $1-K*$ on this interval has norm $O(1+B_c)$ by
Lemma~\ref{lem:brownian-equations}.  Solving the resulting equation for
$d$ proves \eqref{eq:B14}.  To bound its middle term, use
$S*(\delta_0-K)=K_1-K_3$ and write
\[
 (1-K*)_{B_c}^{-1}(K_1-K_3)
    =S_{\mathrm{in}}
      -(1-K*)_{B_c}^{-1}(K*S_{\mathrm{out}}).
\]
The inverse costs a polynomial in $1+B_c$, whereas $S_{\mathrm{out}}$
has an exponential tail.  This gives uniform supremum and $L^1$ bounds;
bounded values of $B_c$ are covered by the same inverse estimate.

Integrating \eqref{eq:B14} shows
\[
 \int\rho_{B_c}\leq b^2/2+O(b).
\]
The total-density asymptotic in \eqref{eq:B5}, together with its relation between
$b$ and $B_*$, yields $B_c\leq B_*+O(1)$.  At $\mathcal W=W$, the
two-sided particle band yields the reverse estimate.  The constants can
be enlarged to include bounded $B_c$.

It remains to check the Jacobian when the single is present.  The pair
block is a diagonally dominant matrix with nonpositive off-diagonal
entries, and its inverse has absolute row sums $O(\mathcal W^{-1})$.
At the symmetric solution the single diagonal entry is
\begin{equation}\label{eq:B15}
 \mathcal W-\sum_j\varphi'(-u_j)
   =2\pi\mathcal W\,
       S*(\rho_{\mathrm{ext}}\mathbf1_{|x|>B_c})(0)+O(1),
 \qquad \rho_{\mathrm{ext}}=1/\pi+K*\rho_{B_c}.
\end{equation}
Here $\rho_{\mathrm{ext}}=\rho_{B_c}$ inside the interval.  To see the
identity, use $\varphi'=2\pi(K_1-K_3)$, apply midpoint summation and
\eqref{eq:B14}, and then convolve
\[
 \rho_{B_c}-K*\rho_{B_c}
   =1/\pi-\rho_{\mathrm{ext}}\mathbf1_{|x|>B_c}
\]
with $S$.  The impurity correction in \eqref{eq:B14} has total mass
$O(\mathcal W^{-1})$, so its contribution is $O(1)$.

The edge lower bound in Proposition~\ref{prop:brownian-asymptotics}
and positivity of $S$ bound the main term in \eqref{eq:B15} below by
\[
 c\mathcal W\sqrt{1+B_c}\,e^{-\pi B_c/2}.
\]
The same conclusion for bounded $B_c$ follows directly from positivity.
The absolute sum of the cross entries is
\[
 \sum_j|\varphi'(u_j)|=O((1+b)\mathcal W).
\]
For example, sum a smooth integrable envelope of $|\varphi'|$ by the
counting rule and use $\rho_{B_c}\leq C(1+B_c)\leq C(1+b)$.
Combining this with the pair inverse bound shows that the correction in
the single Schur complement is $O(1+b)$.  Finally \eqref{eq:B5} and the
circumference assumption give
\[
 W\sqrt{1+B_*}\,e^{-\pi B_*/2}
       \asymp b(\log b)^D.
\]
This dominates the Schur correction and the bounded remainder by a factor
tending to infinity.  The bound improves for $\mathcal W\geq W$, since
$B_c\leq B_*+O(1)$.  The Schur complement is therefore invertible, and
block inversion gives the claimed polynomial bound for the full inverse.
\end{proof}

The center equations are well conditioned except for the explicitly
controlled single-particle direction.  A different estimate is needed
for the imaginary pair spacings: their corrections are exponentially
small, but they are divided by products over all other pairs.  The next
bound controls precisely those products.

\begin{lemma}[Suppression of pair-spacing corrections]
\label{lem:bethe-suppression}
Let the centers be those of Lemma~\ref{lem:bethe-centers}.  Set
\[
 A(x)=\log\frac{|x|\sqrt{x^2+16}}{x^2+4},
 \qquad \kappa(x)=1+(A*\rho_{B_c})(x).
\]
Uniformly for $\mathcal W\geq W$,
\begin{equation}\label{eq:B16}
 \inf_{|x|\leq B_c}\mathcal W\kappa(x)
    \geq c_2\mathcal W\sqrt{1+B_c}\,e^{-\pi B_c/2}
    \gg\log\mathcal W,
\end{equation}
where the last ratio tends to infinity as $b\to\infty$.  Moreover,
\[
 \mathcal W+\sum_{t\ne j}A(u_j-u_t)
    =\mathcal W\kappa(u_j)+O(\log\mathcal W)
\]
uniformly in $j$.
\end{lemma}

\begin{proof}
At the upper member of a formal pair, the external factors in the first
equation of \eqref{eq:B8} contributed by another pair, at center
difference $x$, multiply to
\[
 \frac{x(x+4i)}{x^2+4}.
\]
Its log modulus is $A(x)$.  The Fourier identity
\[
 \widehat A(k)=-\pi\frac{(1-e^{-2|k|})^2}{|k|}
\]
implies
\[
 \kappa=R*(\rho_{\mathrm{ext}}\mathbf1_{|x|>B_c}),
 \qquad \widehat R(k)=\pi\frac{\tanh|k|}{|k|}.
\]
Indeed $\widehat R(1-\widehat K)=-\widehat A$ and $\int R=\pi$.
The partial fraction expansion of $\tanh k/k$, obtained from the product
for $\cosh k$, gives, for $x\ne0$,
\[
 R(x)=2\sum_{m\geq0}
          \frac{e^{-(2m+1)\pi|x|/2}}{2m+1}.
\]
Thus $R$ is positive and $R(x)\geq2e^{-\pi|x|/2}$.  The edge lower bound
for $\rho_{B_c}$ also bounds $\rho_{\mathrm{ext}}$ below by a constant
times $\sqrt{1+B_c}$ in a fixed interval just outside each endpoint.
Convolution over these intervals proves the first inequality in
\eqref{eq:B16}.  Its final comparison follows from the scale estimate
at the end of Lemma~\ref{lem:bethe-centers}; for larger $\mathcal W$ the
left side only gains the needed margin.

For the sum, apply midpoint summation in the coordinate $s=Z_c(u)$.
Away from the cell containing $u_j$ and its two neighbors, monotone or
bounded-variation summation for the logarithmic singularity costs
$O(\log\mathcal W)$.  The spacings from
Lemma~\ref{lem:bethe-centers} show that the log integral over each of
the excluded quantile cells is at most
$C\log\mathcal W/\mathcal W$.  This gives the same total cost for those
cells.  Finally replacing $d$ by $\rho_{B_c}$ uses \eqref{eq:B14}:
the inverse function there is bounded, $A$ is integrable, and the
remaining error is polynomial in $b$ times $\mathcal W^{-1}$.
\end{proof}

\begin{lemma}[Finite-circle roots]\label{lem:bethe-roots}
For the parameter range of this Section, the center solution of
\eqref{eq:B13} has a locally unique continuation to solutions of
\eqref{eq:B8}, with the singular solutions understood as analytic
continuation limits.  At zero twist write these solutions as
\[
 k_{j+}=a_j+i+X_j,\qquad k_{j-}=a_j-i+Y_j,
 \qquad u_j^{\mathrm{true}}=(k_{j+}+k_{j-})/2.
\]
For every fixed $A_0>0$, the differences between the true variables and
their formal center values are $O_{A_0}(\mathcal W^{-A_0})$, uniformly
for $\mathcal W\geq W$, once $b$ is sufficiently large.  The solutions
are locally analytic in $\mathcal W$ and in twists near zero, and these
local solutions link along $[W,\infty)$.  At fixed $n_p$, their zero-twist
centers satisfy
\begin{equation}\label{eq:B18}
 u_j=\frac{\pi(j-(M+1)/2)}{\mathcal W}
             +O_{n_p}(\mathcal W^{-2})
       \qquad (\mathcal W\to\infty).
\end{equation}
Ordinary nonsingular parameters form a generic set on this continued
branch.  A common twist translates every rapidity by $\zeta/\mathcal W$.
\end{lemma}

\begin{proof}
We first rewrite the equations at the formal pair spacings without
dividing by a vanishing internal factor.  The resulting system has a
polynomially bounded inverse linearization, while its error is smaller
than every inverse power of $\mathcal W$; this will give the true roots
by contraction.

The
$j$th auxiliary equation can be written
\[
 \frac{X_j}{2i+X_j}\frac{Y_j-2i}{Y_j}=A_j,
\]
where $A_j$ is analytic and contains only external factors and the twist.
Set $X_j=D_jY_j$ and impose
\[
 D_j=A_j\frac{2i+X_j}{Y_j-2i}.
\]
At the formal solution $|A_j|=1$ and $D_j=-A_j$, so $D_j$ stays away
from zero.  Denote
by $F_j$ the product in the upper-member momentum equation after removing
its internal pair factors.  Solving this equation for $Y_j$ gives the
regular relation
\begin{equation}\label{eq:B17}
 Y_j=(D_j-1)
 \frac{e^{ik_{j+}\mathcal W}(2i+X_j)^2}
      {F_j(4i+X_j-Y_j)D_j^2}.
\end{equation}
In addition, retain the product of the two momentum equations, replacing
the internal auxiliary ratio by $A_j$, and retain the single equation
when present.  The logarithms of these last equations are taken near
\eqref{eq:B13}.  They are analytic there.  Distinct pair centers are
separated; the single need not be separated from a center, because its
potentially dangerous factors have nonzero imaginary shifts.

We give the size estimates for this regular system.  On a neighborhood
of radius $\mathcal W^{-C}$, with a fixed sufficiently large $C$, the
logarithmic derivatives through two orders of the nonzero external
factors, and the row sums of the logarithmic derivatives of their
products, are bounded by fixed powers of $\mathcal W$.
This follows from the center spacing
bound and the polynomial number of factors.  Increasing $C$ therefore
makes the variation of each product logarithm arbitrarily small.
We do not require a polynomial bound for the raw product $F_j$ or its
derivatives.  For $A_j$, whose formal modulus is one, these logarithmic
bounds give polynomial bounds for the ordinary derivatives.  For the
upper-member equation the bounded quantity is instead the normalized
ratio $e^{ik_{j+}\mathcal W}/F_j$.
Indeed, at the formal centers the logarithm of the modulus of this
ratio is
\[
 \log\left|\frac{e^{ik_{j+}\mathcal W}}{F_j}\right|
 =-\mathcal W-\sum_{t\ne j}A(u_j-u_t)+O(\log\mathcal W)
 =-\mathcal W\kappa(u_j)+O(\log\mathcal W).
\]
The possible single-particle factor accounts for the first remainder;
without a single it is absent.  By \eqref{eq:B16}, the negative term
dominates every fixed multiple of $\log\mathcal W$.
The preceding logarithmic derivative bounds preserve this estimate on
a slightly smaller neighborhood.  Hence the fraction after $D_j-1$
in \eqref{eq:B17}, together with the derivatives needed for contraction,
is smaller than every inverse power of $\mathcal W$.

If the right side of \eqref{eq:B17} is set to zero, then $Y=0$.
After eliminating $X=DY$, the linearization is triangular: first the
$Y$ equations, then the center equations \eqref{eq:B13}, and finally
the equations determining $D$.  Its inverse is bounded by a power of
$\mathcal W$, by Lemma~\ref{lem:bethe-centers}.  Apply this inverse to
the regular system and use its second-derivative row-sum bounds.  On
a smaller polynomial-radius ball, the resulting Newton correction map
is a contraction: its forcing term is superpolynomially small and its
Lipschitz constant is less than one.  This proves local existence,
uniqueness, and the claimed error.  The same construction for complex
parameters proves local analyticity.  Its balls can be chosen to
overlap as $\mathcal W$ varies, because the displacement from the
formal solution is much smaller than their radii.  Thus the local
solutions link along the whole ray.

For fixed $n_p$, expansion of \eqref{eq:B13} at large $\mathcal W$
gives \eqref{eq:B18}; the fraction in \eqref{eq:B17} is then
$e^{-\mathcal W}$ times a fixed-$n_p$ polynomial.  The only remaining
issue is whether an internal singularity could persist identically
and prevent use of ordinary Bethe vectors.  At formal spacings,
the external pair factors cancel in $A_j$.  Its explicit twist factor
is $e^{-2i\chi}$; the remaining factor is independent of the other
pair centers and has bounded derivatives, including the single factor
if present.  At this fixed-$n_p$ end, the center and single responses
to the twists are $O_{n_p}(\mathcal W^{-1})$, since the eliminated pair
equations have the fused phases of \eqref{eq:B13}.  Hence an infinitesimal
spin twist makes $D_j-1$ nonzero generically even if it vanishes at zero
twist.  The exponentially small correction in \eqref{eq:B17} does not
change this conclusion.  Analytic continuation gives generic ordinary
points along the branch.  Finally direct substitution in
\eqref{eq:B8} proves the common translation by $\zeta/\mathcal W$.
At real zero twist, conjugation symmetry, or continuity from the
ordinary roots, gives real eigenenergies whenever the vectors are
nonzero; no global uniqueness assertion about Bethe roots is needed.
\end{proof}

\subsection{Identification with the simple ground state}

The construction has now produced roots at every required circumference.
To use them in the transfer norm, we must still show that their plane
waves represent the balanced ground state.  We first examine their
coefficients at the fixed-$n_p$ end $\mathcal W\to\infty$, and then use
simplicity to continue the identification back to $W$.

\begin{proposition}[The ground branch]\label{prop:bethe-ground}
At zero twist, the continuation-limit plane wave associated with
Lemma~\ref{lem:bethe-roots} has energy equal to the balanced ground
energy of $H_2$ on the circle of circumference $W$.  At generic parameters
the plane wave is nonzero; at exceptional parameters the energy and
eigenstate are understood by continuation.  Thus its transfer eigenvalue
is $\|T_{n_p}\|$.  Infinitesimal twisted eigenvalues near this simple
ground level may be used in its computation.
\end{proposition}

\begin{proof}
We first show that the algebraic wave is nonzero and find its limit as
$\mathcal W\to\infty$ at fixed $n_p$.  The positive-overlap criterion
will then identify it at that end; simplicity will keep it on the ground
level as $\mathcal W$ decreases to $W$.

Order the initial momenta in adjacent pairs, lower member before upper
member, with the single last if present.  We first determine the
coefficient produced by the $B$ operators of Lemma~\ref{lem:bethe-algebra}.
At formal pair spacings, the spin-$\frac12$ product on a pair, evaluated
at its center, maps into the singlet.  Its two factors are proportional
to an antisymmetric projection and a permutation, which proves this
directly.  At an arbitrary spectral value the pair product preserves
that singlet and acts by the scalar
\[
 (u-u_j)^2+4.
\]
Consequently, applying $B(u_j)$ lowers the designated pair into its
singlet while acting on previously produced singlets by these scalars.
The vacuum-preserving factors on the other pairs can vanish only when
pair centers coincide.  The real single, if present, remains separated
from their dangerous factors by imaginary shifts.  Hence the initial
coefficient is nonzero and, after division by a scalar, tends to the
product of pair singlets in each index.  At the spacings
\eqref{eq:B18}, this scalar and its reciprocal have fixed-$n_p$
polynomial bounds.  Replacing the formal roots by the true roots changes
this normalized coefficient by an exponentially small quantity times
such a polynomial.

We must also control coefficients for permutations that reverse a pair.
By braid consistency, each can be written as a $B$ construction with
permuted physical sites, multiplied by scalar exchanges on the joint
all-up vacuum.  Reversing the internal order can cost a factor with
denominator $|X_j-Y_j|$.  On the other hand, each single-index $B$
vector vanishes at the exact string variables when that pair alone is
reversed.  This can be seen without exchanging its sites.  On the
all-up vector, the coefficient of the vector with its only down spin
at site $r$ in $B(a)$ is
\[
 -2i\prod_{s>r}(a-k_s-i)\prod_{s<r}(a-k_s+i).
\]
Suppose the upper member $k_p=a+i$ precedes the lower member
$k_q=a-i$, so $p<q$.  If $r<q$, the first product vanishes at $s=q$;
if $r\geq q$, the second product vanishes at $s=p$.  Thus $B(a)$
annihilates the vacuum for any separation of these sites and any values
of the other momenta.  Since the $B$ entries commute, apply this $B(a)$
first in the full spin vector.  This proves the claimed vanishing as
a polynomial identity, without a potentially singular conjugation.

Apply this vanishing separately to each reversed pair.  Taylor's
formula in the corresponding variables bounds each single-index
vector by a product of factors $|X_j|+|Y_j|$.
Precisely, for each reversed pair subtract the evaluation at
$X_j=Y_j=0$.  These evaluation operators commute because their variable
pairs are disjoint, and each subtracted evaluation is zero by the
preceding identity, with all other variables unrestricted.  The product
of the resulting difference operators is the original vector.  Applying
the integral form of Taylor's formula to this product gives the claimed
product bound, with fixed-$n_p$ polynomial constants at
\eqref{eq:B18}.  The two indices supply
the square of this product, whereas the vacuum scalar costs only one
factor $|X_j-Y_j|^{-1}$ per reversal.  Since
\[
 X_j-Y_j=(D_j-1)Y_j,
 \qquad
 Y_j=(D_j-1)\times\hbox{the fraction in \eqref{eq:B17}},
\]
the quotient is removable even at $D_j=1$.  At fixed $n_p$ its bound is
\[
 \operatorname{poly}_{n_p}(\mathcal W)e^{-r\mathcal W}
\]
when $r$ pairs are reversed.  All remaining permutation factors cost
only fixed-$n_p$ powers, because distinct centers are separated.
The same argument works with infinitesimal twists and then passes
to zero twist.  In particular the initial ordered amplitude remains
nonzero at the large-circle end, and singular roots there yield ordinary
continuation-limit eigenfunctions.

To apply Corollary~\ref{cor:bethe-ground-test}, rescale the circle of
length $\mathcal W$ at coupling $2$ to a fixed length $W_0$.
The new coupling is $c=2\mathcal W/W_0$, momenta are multiplied by
$\mathcal W/W_0$, and energies by $(\mathcal W/W_0)^2$.
The eigenenergy is
\[
 -Mc^2/2+O_{n_p}(1).
\]
Consider a coordinate chart with the seam avoided, distinct pair centers
separated by a fixed positive distance, and within-pair gaps rescaled
by $c$.  Reversed-pair amplitudes are negligible there, because the
spans remain short of a full circumference.  A permutation that pairs
positions near different separated centers is negligible for the same
exponential reason.  The remaining permutations exchange entire blocks.
Their coefficients are given, to leading order, by the two fusion
scalars in the line calculation.  As $c\to\infty$, the pair--pair scalar
tends to $-1$, and the pair--single scalar tends to $1$.

The momentum conjugate to a pair center is twice its individual momentum.
Thus \eqref{eq:B18} gives the center wave
\[
 \det\left[
   \exp\!\left\{\frac{2\pi i}{W_0}
        \left(j-\frac{M+1}{2}\right)q_a\right\}
     \right]_{j,a=1}^M,
\]
together with the normal profile \eqref{eq:B12} and a constant single
wave if present.  For ordered distinct centers, the Vandermonde formula
is a constant phase times
$\prod_{a<t}\sin(\pi(q_t-q_a)/W_0)$.  This product is strictly positive.
After the static sign change, each double singlet has nonnegative components,
as does the possible majority-color single.

The same phase choice applies to charts crossing the seam, by periodicity
and translation.  To check the latter point, sum \eqref{eq:B13} to get
$2\sum_j u_j+\ell v=0$.  At zero twist, periodicity quantizes the total
momentum of a nonzero true wave in $(2\pi/\mathcal W)\mathbb Z$.
Its difference from the formal value is $o(\mathcal W^{-1})$ at the
fixed-$n_p$ large-circle end, so that integer is zero there.  The integer
is constant along the continuous root branch, including continuation
limits, and hence the total momentum is zero throughout.

The convergence just obtained is uniform on compact separated-center
charts with bounded rescaled gaps, before the final $L^2$ normalization.
After normalization, the wave on each such chart is a common scalar
multiple of this limiting shape, up to a vanishing relative error.
The scalar is bounded above because the shape has positive norm on
any one nonempty chart.  It cannot tend to zero along a subsequence:
Lemma~\ref{lem:bethe-pair-compactness} would then give zero limit on
every retained chart, although those charts retain the full norm.
Every normalized limit consequently has the same strictly positive
center shape, with the chart cutoffs included.
Corollary~\ref{cor:bethe-ground-test} identifies the plane branch with
the ground state for sufficiently large $c$.

Finally decrease $\mathcal W$ along the branch of
Lemma~\ref{lem:bethe-roots}.  The energy is locally analytic and is an
eigenvalue at generic parameters.  The plane coefficients and $B$
vectors are meromorphic; their nonvanishing at the large-circle end
shows that ordinary nonzero states occur on an open dense set.  At
complementary parameters, take the twist continuation first and then
pass to zero twist.  Continuity of the compact-resolvent spectrum on
a finite circle keeps the limiting energy in the spectrum.  The
balanced ground energy remains simple by
Lemma~\ref{lem:brownian-ground}, so a continuous eigenvalue branch
already equal to it cannot leave it without a forbidden degeneracy.
The identification therefore holds all the way to $\mathcal W=W$.
Commutation and positivity in Lemma~\ref{lem:brownian-ground} now
identify the corresponding transfer eigenvalue with $\|T_{n_p}\|$.
The same simple eigenvalue admits the infinitesimal twist continuation
used below.
\end{proof}

\subsection{The transfer eigenvalue on a finite circle}
\label{subsec:bethe-transfer}

We now set $c=2$ and $\mathcal W=W$.  Proposition~\ref{prop:bethe-ground}
identifies the solution of \eqref{eq:B8} on which the transfer norm is
attained.  We must still calculate that eigenvalue for the heat kernel on
the circle.  The rotor representation carried across the seam contributes
a sum which is absent from a calculation with infinite-line plane waves.
We will first obtain this sum exactly and then bound its nontrivial terms
uniformly when $n_p$ grows with $W$.

Write $n_p=2M+\ell$, with $\ell\in\{0,1\}$, and retain the twists
$\zeta,\chi$ used in \eqref{eq:B8}.  We conjugate each exterior binary
index by the antisymmetric tensor.  This fixed change of basis reverses
the diagonal matrix-index twists.  The seam therefore acts in rotor
sector $(q,j)$ by the scalar $e^{i\zeta q}$ and by $g_j^{-1}$ on each of
its two matrix indices, where $g_j$ is the spin-$j$ representation of
$g=\operatorname{diag}(e^{i\chi},e^{-i\chi})$.  Throughout this calculation
\[
 q\in\mathbb Z,\qquad j\in\tfrac12\mathbb Z_{\ge0},\qquad
 2j\equiv q\pmod 2,\qquad d_j=2j+1,\qquad
 E_{qj}=\tfrac12(q^2+d_j^2-1).
\]
The normalized Haar harmonics are $\sqrt{d_j}\,D^j_{ab}(Q)e^{iqv}$.
By Lemma~\ref{lem:brownian-ground}, the simple balanced ground space of
\eqref{eq:B2} is invariant under the transfer operator.  Near zero twist
its analytic continuation is still one-dimensional.  Denote the transfer
scalar on that space by $\Lambda_{n_p}(\zeta,\chi)$; at zero twist it is
$\|T_{n_p}\|$.

\begin{lemma}[Finite-circle transfer formula]
\label{lem:bethe-transfer-formula}
For the root branch of Lemma~\ref{lem:bethe-roots}, with fixed finite
$n_p$ and $W$, let $k_t$ be its momenta and let $t_j$ be the auxiliary
trace eigenvalue in \eqref{eq:B9}.  The function
$\Lambda_{n_p}(\zeta,\chi)$ is given near zero twist by
\begin{equation}
 \Lambda_{n_p}(\zeta,\chi)
 =\lambda^{n_p}\sum_{q,j}e^{-WE_{qj}+i\zeta q}
 \frac{(-2)^{n_p}i^{-2n_p}t_j(-iq)^2}
 {\displaystyle\prod_{t=1}^{n_p}
   \bigl[(d_j+1)^2-(ik_t-q)^2\bigr]
   \bigl[(d_j-1)^2-(ik_t-q)^2\bigr]} .
 \label{eq:B20}
\end{equation}
The sum is over the admissible rotor sectors just specified.  Apparent
singularities in this expression are interpreted by analytic
continuation from parameters at which coefficient extraction and the
Bethe vectors are nonsingular.
\end{lemma}

\begin{proof}
At an ordinary point of the root branch, choose a momentum ordering
whose color coefficient $A$ in the Bethe wave is nonzero.
The transfer preserves its eigenline, so the coefficient of
$\exp(i\sum_a k_aY_a)$ in the transferred wave is
$\Lambda_{n_p}A$.  We compute this coefficient, initially at the
nonresonant parameters specified below.

We first compute the contribution near one output coordinate $Y$.  Let
$X$ be the integration coordinate assigned to it, let $k$ be the momentum
carried by $X$, and fix the rotor sector $(q,j)$ on the intervals outside
these two insertions.  Put $d=2j+1$ and $z=ik-q$.  If $S_j$ denotes the
spin-$j$ angular momentum, the operators
$\eta_L=S_j\cdot\sigma$ and $\eta_R=S_j\cdot\sigma$ act on the left and
right rotor and binary indices, respectively.  These two operators
commute.  We extract the coefficient of $e^{ikY}$ obtained by taking
$Y$ as the integration endpoint for $X$.  The sum of this coefficient
over the two orders $X<Y$ and $X>Y$ is
\begin{equation}
 \frac{-2(z-2\eta_L)(z-2\eta_R)}
 {\bigl((d+1)^2-z^2\bigr)\bigl((d-1)^2-z^2\bigr)}.
 \label{eq:B19}
\end{equation}
We verify both the normalization and the endpoint denominators in this
formula.

Suppose first that $X<Y$, with no other insertion between them.
Multiplication by the conjugate entry at $X$ changes the intermediate
sector to $(q-1,j\pm\tfrac12)$.  Its energy minus $E_{qj}$ is
$1\pm d-q$.  On one matrix index the Clebsch--Gordan projections in
$j\otimes\tfrac12$ are
\[
 P_+(\eta)=\frac{j+1+\eta}{d},\qquad
 P_-(\eta)=\frac{j-\eta}{d}.
\]
The two matrix indices supply $P_\pm(\eta_L)P_\pm(\eta_R)$, with the same
choice of sign because they belong to the same intermediate rotor
representation.  The normalization of the Haar harmonics supplies the
additional factor $d/(d\pm1)$.  Indeed, multiplication from a normalized
spin-$j$ matrix element to a normalized spin-$(j\pm\tfrac12)$ matrix
element has coefficient $\sqrt{d/(d\pm1)}$ times the two
Clebsch--Gordan coefficients.  The return insertion is its adjoint and
contributes the same factor.  Thus this
dimension factor occurs once in the product of the two insertions, not
once for each matrix index.

For $X<Y$, the input plane times the heat factor is
$e^{ikX}e^{-(1\pm d-q)(Y-X)}$.
Its upper-endpoint contribution is $e^{ikY}/(1\pm d+z)$.
For $X>Y$ the lower-endpoint calculation gives denominator
$1\pm d-z$.  The exterior transpose and the conjugation by the
antisymmetric tensor replace $\eta_L,\eta_R$ by $-\eta_L,-\eta_R$ in
that contribution.  Consequently the full coefficient is
\[
 \sum_{s=\pm1}\frac{d}{d+s}
 \left\{
 \frac{P_s(\eta_L)P_s(\eta_R)}{1+sd+z}
 +\frac{P_s(-\eta_L)P_s(-\eta_R)}{1+sd-z}
 \right\}.
\]
Substituting the two projection formulas and taking a common denominator
gives \eqref{eq:B19}.  If $d=1$, the minus representation is absent:
we omit it before simplifying.  The result is the canceled value of
\eqref{eq:B19}, in particular $2/(4+k^2)$ when $(q,j)=(0,0)$.

We next extract the coefficient of a full output plane
$\exp(i\sum_a k_aY_a)$ in the actual heat-kernel integral.
Initially choose a nonsingular Bethe vector and parameters for which
all integrated partial-sum rates are nonzero, no $ik_a$ is a pure
rotor-energy difference, and no $i(k_a-k_t)$, for $a\ne t$, is such a
difference.  The relevant energy differences are even integers.
These conditions can be met on the branch: a generic common twist
$\zeta$ removes pure-rotor coincidences and zero partial-sum rates.
Near the large-circle end, the separated centers in \eqref{eq:B18}
and an infinitesimal spin twist $\chi$ resolve the remaining
pair-difference resonances.  Here $n_p$ and $c=2$ are fixed while the
circumference varies.  We will pass to all other parameters by analytic
continuation after computing the coefficient.

Fix a total ordering of $X$ and $Y$ and expand in rotor harmonics.
An ordered block $x_1<\cdots<x_r$ of input coordinates in a gap
$(a,b)$ between consecutive output coordinates has exponential rates
$\alpha_1,\ldots,\alpha_r$.  Repeated integration produces terms with
endpoint phases
\[
 \exp\left\{a\sum_{h=1}^{s}\alpha_h
              +b\sum_{h=s+1}^{r}\alpha_h\right\},
 \qquad 0\le s\le r.
\]
Thus a prefix is assigned to the left endpoint and the remaining
suffix to the right; this description includes the nested integrations
whose intermediate endpoints coincide.  Over all gaps it gives a
weakly order-preserving assignment of inputs to outputs or to the
seam endpoints $0,W$.

At an output $Y_a$, the exponent is the sum of the assigned input
momenta, multiplied by $i$, plus the adjacent rotor-energy difference.
If no input is assigned there, the exponent is purely a rotor-energy
difference and cannot equal $ik_a$.  Every output therefore receives
at least one input.  Since their numbers agree, exactly one input is
assigned to each output and none to the seam.  The difference
nonresonance then forces that input to carry $k_a$, and forces its
remaining rotor-energy difference to vanish.

In particular, a contributing gap contains at most two inputs, one
assigned to each endpoint.  When two integration variables lie in the
same $Y$ gap and are assigned
to its opposite ends, the coefficient is the product of their two
endpoint factors.  This can be checked without separating the ordered
integrations: for exponential rates $\alpha,\beta$, the term assigning
$x_1$ to $a$ and $x_2$ to $b$ in
\[
 \int_a^b e^{\beta x_2}\int_a^{x_2}e^{\alpha x_1}\,dx_1\,dx_2
\]
has coefficient $-1/(\alpha\beta)$, precisely the lower and upper
endpoint coefficients multiplied together.  The other terms assign the
two variables to a common endpoint and do not contribute to the required
plane.  After every $X$ has been paired with a $Y$, the vanishing
energy differences show that all intervening intervals carry the
same rotor energy.  The charge has already returned to the same $q$,
so equality of energies forces the same $j$.
Thus one fixed through sector $(q,j)$ contributes the ordered product of
\eqref{eq:B19} over the momenta, followed by its seam trace.

To justify coefficient extraction also for the infinite rotor sum,
work in the ordered complex tube
$0<\operatorname{Re}Y_1<\cdots<\operatorname{Re}Y_{n_p}<W$.
At fixed finite $n_p$, intermediate labels differ from the through
labels by bounded steps, so their energies differ from $E_{qj}$ by
$O_{n_p}(1+|q|+d_j)$.  On each compact subset of this tube, the
endpoint terms are therefore bounded by a fixed polynomial in the
labels times
\[
 \exp\{-WE_{qj}+C(1+|q|+d_j)\}.
\]
The quadratic energy gives normal convergence of the analytic
exponential sums, including after any finite imaginary translation.
Modulo even-integer shifts, their exponent vectors have only finitely
many types.  The commuting translations $Y_a\mapsto Y_a+i\pi$
multiply an exponential of rate $\gamma_a$ by $e^{i\pi\gamma_a}$;
these multipliers distinguish the types.  Polynomial interpolation in
the translations consequently projects onto each type.  In the type
containing the required plane, divide by $\exp(i\sum_a k_aY_a)$.
The resulting locally convergent Laurent series in $w_a=e^{2Y_a}$,
on $1<|w_1|<\cdots<|w_{n_p}|<e^{2W}$, has the desired plane
coefficient as its constant coefficient.  Uniqueness
of Laurent coefficients proves that no other term of the infinite
sum changes the extraction above.

Finally,
\[
 i(z-2\eta)=(-iq-k)-2iS_j\cdot\sigma
           =\mathcal L_j(-iq-k).
\]
The ordered product in each matrix index is therefore the monodromy
whose twisted trace has eigenvalue $t_j(-iq)$ in \eqref{eq:B9}.
The two indices give its square.  Multiplying by $\lambda^{n_p}$ and
summing the seam weights yields \eqref{eq:B20}.  At zero twist this is
the top transfer eigenvalue by Proposition~\ref{prop:bethe-ground} and
Lemma~\ref{lem:brownian-ground}; exceptional ring lengths follow by
continuity.
\end{proof}

\subsection{Uniform control of the nontrivial rotor sectors}
\label{subsec:bethe-winding}

We now estimate \eqref{eq:B20} in the regime of
Proposition~\ref{prop:brownian-parity}.  Its $(q,j)=(0,0)$ term is
\[
 Z_{n_p}(\zeta,\chi)
   :=\prod_{t=1}^{n_p}\frac{2\lambda}{4+k_t(\zeta,\chi)^2}.
\]
The issue is to compare every other term with this product without
accumulating a fixed loss for each of the $O(Wb^2)$ particles.

\begin{lemma}[Uniform transfer asymptotics]
\label{lem:bethe-transfer-uniform}
Fix $A,D>0$ and positive constants $c_-,c_+$.  Suppose $b\to\infty$,
\[
 c_-e^{\pi b}b^{-1/2}(\log b)^D
 \le W\le c_+e^{\pi b}b^{-1/2}(\log b)^D,
 \qquad |n_p/W-b^2|\le Ab.
\]
There is $c'>0$, depending only on these fixed constants, such that at
zero twist the sum in \eqref{eq:B20} over $(q,j)\ne(0,0)$, divided by
$Z_{n_p}(0,0)$, is $O(e^{-c'W})$, uniformly over these particle numbers.
If $u_r$ are the real centers of \eqref{eq:B13} at $\mathcal W=W$, then
\begin{equation}
 \log\|T_{n_p}\|
   =\sum_{r=1}^M p(u_r)+\ell f(0)+o(W^{-1}),
 \label{eq:B21}
\end{equation}
with the error uniform over the same band.
\end{lemma}

\begin{proof}
Set the spin twist to zero, initially retaining continuation from
arbitrarily small nonsingular spin twists where needed.  In
\eqref{eq:B9}, evaluated at $-iq$, every shifted argument of $Q_0$
has an odd integral imaginary part relative to a real formal root.
Indeed $2m\equiv2j\equiv q\pmod2$, so $q+2m\pm1$ is odd, and the
same is true of $q\pm(2j+1)$.  These factors are uniformly separated
from zero.

At the formal pair spacing $k_{r\pm}=u_r\pm i$, the contribution of
this pair to $F$ and its accompanying $Q_0$ factors in each summand of
\eqref{eq:B9} cancels to
\[
 (-iq-u_r)^2+d^2,\qquad d=2j+1,
\]
independently of $m$.  It follows that, in every one of the $d^2$
summands obtained by squaring $t_j$, the absolute contribution of this
pair relative to its $(0,0)$ contribution is
\[
 \frac{(u_r^2+1)(u_r^2+9)}
 {\left|\bigl((iu_r-q)^2-(d+2)^2\bigr)
              \bigl((iu_r-q)^2-(d-2)^2\bigr)\right|}
 \le1.
\]
To see the inequality, factor the denominator into four linear factors.
Their absolute values are $\sqrt{u_r^2+w^2}$, where the four widths are
$|q\pm(d+2)|$ and $|q\pm(d-2)|$.  All are positive odd integers, and
at least two are at least three.  Their product is consequently at
least $(u_r^2+1)(u_r^2+9)$.

This comparison remains uniform for the true roots.  The linear factors
just used, and the individual factors before cancellation for the two
pair constituents, stay a fixed distance from zero.  Equations
\eqref{eq:B16}--\eqref{eq:B17} make the root corrections smaller than
every power of $W^{-1}$.  There are only polynomially many factors in
$W$ in the particle band.  Hence the total multiplicative change is
bounded, and in fact differs from one by less than every prescribed
power of $W^{-1}$.  This is the point at which an estimate per particle
would have been insufficient without the sign-free bound of one at the
formal spacing.

If $\ell=1$, denominators contributed by the single momentum can vanish
at zero common twist before cancellation.  We control them on one
complex $\zeta$ circle instead.  Choose a sufficiently large fixed
$C$ and put $|\zeta|=W^{-C}$.  The common twist translates every
momentum and auxiliary root by $\zeta/W$.  At $\zeta=0$ the single
momentum is zero up to an error smaller than every power of $W^{-1}$.
Thus on this circle every possibly vanishing single factor has modulus
bounded below by a fixed multiple of $W^{-C-1}$.  All single-root
contributions are consequently bounded by a fixed power of
$1+W+d+|q|$.  The change in all pair factors together is bounded as
well: each argument moves by $W^{-C-1}$ and the number of factors is
polynomial in $W$.  Choosing $C$ once sufficiently large controls their
product.

We now justify using the maximum-modulus principle on this circle;
this requires analyticity on its whole interior, with no appeal to a
lower bound for the gap between different total-momentum sectors.
Let $\psi_0$ be the normalized untwisted ground vector, with total
momentum zero, and define the change of gauge
\[
 (U_\zeta\psi)(X)
 =\exp\!\left(\frac{i\zeta}{W}\sum_{t=1}^{n_p}X_t\right)\psi(X).
\]
After this change of gauge, the Hamiltonian on periodic functions is
$H_2+2\zeta P/W+n_p\zeta^2/W^2$, where
$P=-i\sum_t\partial_{X_t}$.  In the $P=0$ subspace its ground line is
therefore the same simple line for every $\zeta$.
Write $T_{n_p}(\zeta,0)$ for the twisted transfer.  Its rotor expansion
is analytic, and so is the fixed-gauge scalar
\[
 \left\langle\psi_0,
 U_\zeta^{-1}T_{n_p}(\zeta,0)U_\zeta\psi_0\right\rangle.
\]
Near zero this scalar equals $\Lambda_{n_p}(\zeta,0)$ by commutation.
It therefore continues $\Lambda_{n_p}$ across any crossings with
other momentum sectors.  The coefficient formula may first be used
with arbitrarily small spin twists and then continued to zero spin
twist, equivalently taking the removable analytic limits in
\eqref{eq:B20}.

On the $\zeta$ circle, the ratio of the nontrivial-sector sum to
$Z_{n_p}$ is therefore bounded by
\[
 C\sum_{(q,j)\ne(0,0)}
 d^2(1+W+d+|q|)^K
 \exp\{-WE_{qj}+W^{-C}|q|\}
 =O(e^{-c'W})
\]
for fixed $K$ and some $c'>0$.  Every nontrivial admissible sector has
$E_{qj}\ge2$, and the energies grow quadratically in $q,d$, so the
polynomial factors and the number of summands are absorbed by the
exponential.  The ratio is analytic on the enclosed disk:
$Z_{n_p}$ is analytic and nonzero there, and the whole nontrivial-sector
sum is $\Lambda_{n_p}-Z_{n_p}$.  The maximum-modulus principle gives
the same bound at zero twist.  For $\ell=0$ the separated pair factors
already give the bound directly; the same circle argument also applies.

It remains to take the logarithm of the leading product.  Since
$2\lambda=b^2$, a formal pair contributes
\[
 \frac{b^4}{(4+(u_r+i)^2)(4+(u_r-i)^2)}
 =\frac{b^4}{(u_r^2+1)(u_r^2+9)}=e^{p(u_r)},
\]
and the possible zero single contributes $b^2/4=e^{f(0)}$.
The cumulative error from replacing the true roots by the centers of
\eqref{eq:B13} is smaller than every power of $W^{-1}$ by
\eqref{eq:B17}; the nontrivial-sector error is exponentially smaller.
The positive zero-twist eigenvalue therefore has the logarithm
\eqref{eq:B21}.
\end{proof}

\subsection{Maximizing over the two particle parities}
\label{subsec:bethe-parity-max}

The transfer asymptotics have reduced the problem to a sum over real
pair centers.  The dressed energy $\epsilon$ of \eqref{eq:B3}, with
endpoint $B=B_*$, expresses this sum as a common bulk term, an explicit
single-particle cost, and a quadrature error.  Recall that
$\epsilon_{\rm in}=\epsilon\mathbf1_{(-B_*,B_*)}$ is also the positive
part of $\epsilon$.

Let $d(x)$ now denote the globally defined counting derivative for
\eqref{eq:B13} at $\mathcal W=W$; it is not truncated to its counting
interval $(-B_c,B_c)$.  Summing $p=\epsilon-K*\epsilon_{\rm in}$ at
the centers and using the definition of $d$ gives the exact identity
\begin{equation}
 \sum_{r=1}^M p(u_r)+\ell f(0)
 =\frac{W}{\pi}\int_{\mathbb R}\epsilon_{\rm in}(x)\,dx-\ell m_*
  +\left[\sum_{r=1}^M\epsilon(u_r)
        -W\int_{\mathbb R}\epsilon_{\rm in}(x)d(x)\,dx\right].
 \label{eq:B22}
\end{equation}
Indeed,
\[
 W\int\epsilon_{\rm in}d
 =\frac{W}{\pi}\int\epsilon_{\rm in}
   +\sum_{r=1}^M(K*\epsilon_{\rm in})(u_r)
   -\ell\bigl((K_1-K_3)*\epsilon_{\rm in}\bigr)(0),
\]
and the remaining single term is $-\ell m_*$ by its definition.

\begin{lemma}[The parity variational error]
\label{lem:bethe-parity-error}
Under the parameter assumptions of Lemma~\ref{lem:bethe-transfer-uniform},
the bracket in \eqref{eq:B22} is at most $C/W$, uniformly in the particle
band.  In each parity there is a particle number in a band of the same
form for which the bracket is at least $-C/W$.  Specifically, one may
take $M$ to be a nearest integer to $W\int\rho_{B_*}$ and set
$n_p=2M+\ell$.
\end{lemma}

\begin{proof}
For the upper bound, replace each $\epsilon(u_r)$ by its positive part.
Use the counting coordinate
\[
 s(x)=\int_0^x d(y)\,dy,
 \qquad s(u_r)=\frac{r-(M+1)/2}{W},
 \qquad s(\pm B_c)=\pm\frac{M}{2W}.
\]
The centers are exactly the midpoints of cells of length $1/W$ in this
coordinate.  The estimates accompanying \eqref{eq:B14} give $d\ge c>0$
and polynomial bounds in $b$ for its fixed-order derivatives on the
needed interval.  The integral equation \eqref{eq:B3} gives such bounds
for the fixed-order derivatives of $\epsilon$.  Moreover
$|\epsilon'|\le\pi/2$ on the whole line by
Proposition~\ref{prop:brownian-asymptotics}.  Thus, for
the inverse function $x=x(s)$,
\[
 \left|\frac{d}{ds}\epsilon(x(s))\right|
 =\left|\frac{\epsilon'(x(s))}{d(x(s))}\right|\le C,
\]
while each further fixed-order derivative is polynomially bounded in
$b$.

Set $g(s)=\epsilon_+(x(s))$.  On any consecutive collection of full
cells in a smooth portion, with outer endpoints $\alpha,\beta$, the
composite midpoint expansion gives
\[
 \sum_{\text{these cells}}g(s_r)
       -W\int_\alpha^\beta g(s)\,ds
 =-\frac{g'(\beta)-g'(\alpha)}{24W}
       +O\!\left(\frac{\operatorname{poly}(b)}{W^2}\right).
\]
The first derivatives are uniformly bounded, and the higher derivatives
and interval length cost only polynomial factors in the remainder.
Thus this error is $O(W^{-1})$ with a constant independent of $b$.
At most two cells meet the cutoffs
$x=\pm B_*$.  Since $\epsilon$ vanishes there and its derivative in
$s$ is bounded, these cells contribute $O(1/W)$ as well.  Consequently
\[
 \sum_r\epsilon(u_r)
 \le W\int_{-B_c}^{B_c}\epsilon_+(x)d(x)\,dx+\frac{C}{W}
 \le W\int_{\mathbb R}\epsilon_{\rm in}(x)d(x)\,dx+\frac{C}{W}.
\]
The last inequality uses positivity: any positive portion outside the
counting interval is only omitted on its left-hand side.

For the lower bound choose $M$ as in the statement.  Equation
\eqref{eq:B5} and the relation between $b$ and $B_*$ imply
$2\int\rho_{B_*}=b^2+O(b)$, so these particle numbers satisfy the
band condition after enlarging its fixed constant.  Since $\int_{-B_c}^{B_c}d=M/W$,
\eqref{eq:B14}, including its uniform $L^1$ bound for the impurity
correction, yields
\[
 \left|\int\rho_{B_c}-\int\rho_{B_*}\right|=O(W^{-1}).
\]
The densities increase with their interval: this follows by positive
iteration of their defining equation.  Their difference in total mass
therefore bounds from below the mass on the excess or missing interval.
First $\rho_B\ge1/\pi$ gives $|B_c-B_*|=O(W^{-1})$.  Those intervals
then lie within distance one of the relevant larger endpoint; applying
the edge bound $\rho_B\ge c\sqrt B$ from \eqref{eq:B5} improves this to
\[
 |B_c-B_*|\le\frac{C}{W\sqrt{B_*}}.
\]
On this excess or missing interval the counting density obeys
$d\le CB_*$, by \eqref{eq:B14} and its defining convolution, and
$|\epsilon(x)|\le C\,\bigl||x|-B_*\bigr|$.  Therefore replacing the
counting interval by $(-B_*,B_*)$ in the dressed-energy integral costs
at most
\[
 CWB_*|B_c-B_*|^2\le\frac{C}{W}.
\]
The same midpoint estimate, now for the smooth function $\epsilon$
on the counting interval, compares its sum and integral with error
$O(1/W)$.  Combining these two bounds proves the lower estimate for the
bracket.
\end{proof}

We can now complete Proposition~\ref{prop:brownian-parity}.  By its
particle-band hypothesis, every maximizing particle number is covered
by the uniform upper bound above.  After one fixed enlargement of the band constant, the choices in
Lemma~\ref{lem:bethe-parity-error} are covered as well and give the
matching lower bounds.  Thus
\[
 \log\tau_\ell
 =\frac{W}{\pi}\int\epsilon_{\rm in}-\ell m_*+O(W^{-1}),
 \qquad \ell=0,1.
\]
By \eqref{eq:B5},
\[
 e^{-\pi B_*/2}=4\pi e B_*e^{-\pi b}(1+o(1)),
 \qquad
 m_*\sim32\sqrt{B_*}\,e^{-\pi b}
       \sim32\sqrt2\,b^{1/2}e^{-\pi b}.
\]
The prescribed size of $W$ gives $Wm_*\asymp(\log b)^D\to\infty$.
Subtracting the two parity formulas gives
\[
 \log\tau_0-\log\tau_1=m_*+O(W^{-1}),
\]
and therefore proves \eqref{eq:B1}.

